\chapter{How to Add Features}
\label{ch:how-to-add-features}

\textit{This chapter provides seven self-contained recipes for extending
\sysname{}.  Each recipe identifies what you are adding, which files to
touch, the concrete steps, a code snippet, a verification command, and
common gotchas.  Work through them in order the first time---Recipe~17.1 is
a two-minute change that builds confidence for the harder ones.}

\begin{keyinsight}
Start with Recipe~17.1 (Add a Canonical Mapping).  It is the smallest
possible change that touches the entity resolution pipeline, and the
benchmark will immediately tell you whether you got it right.
\end{keyinsight}

% ─────────────────────────────────────────────────────────────────────
\section{Add a Canonical Mapping}
\label{sec:add-canonical-mapping}

\noindent\textbf{What you are adding.}
A new alias $\rightarrow$ canonical-name entry in the entity resolution
dictionary so that a surface form like \ic{k8s} resolves to
\ic{Kubernetes}.

\medskip\noindent\textbf{Files to touch.}

\begin{itemize}
    \item \codefile{apps/api/models/ontology.py} --- the
          \ic{CANONICAL\_NAMES} dictionary (currently 534 entries).
\end{itemize}

\medskip\noindent\textbf{Steps.}

\begin{enumerate}
    \item Open \codefile{apps/api/models/ontology.py}.
    \item Locate the \ic{CANONICAL\_NAMES} dictionary (around line~271).
    \item Add your mapping inside the appropriate category section.  Keys
          are lowercase; values are the display-cased canonical form.
    \item Save the file.
    \item Run the synthetic benchmark to verify.
\end{enumerate}

\medskip\noindent\textbf{Code snippet.}

\begin{lstlisting}[style=python, title=Adding a Kubernetes alias]
CANONICAL_NAMES: dict[str, str] = {
    # ... existing entries ...
    # ============================================================
    # CONTAINERIZATION & ORCHESTRATION
    # ============================================================
    "kubernetes": "Kubernetes",
    "k8s": "Kubernetes",       # already present
    "kube": "Kubernetes",      # already present
    "k8": "Kubernetes",        # <-- NEW: common typo/shorthand
    # ...
}
\end{lstlisting}

\medskip\noindent\textbf{Verification.}

\begin{lstlisting}[style=bash, numbers=none]
cd apps/api
python -m evaluation.synthetic_benchmark
\end{lstlisting}

The benchmark generates surface-form variants for every canonical name and
checks that each resolves correctly.  A passing run confirms your new
mapping is well-formed.

\medskip\noindent\textbf{Gotchas.}

\begin{itemize}
    \item Keys \emph{must} be lowercase.  The lookup calls
          \ic{name.lower()} before hitting the dictionary.
    \item Do not map a key that already exists with a different canonical
          value---Python will silently overwrite the earlier entry.
    \item If your canonical name is new (not yet in the dictionary as a
          value), also add the well-known abbreviations to
          \ic{KNOWN\_ABBREVIATIONS} in
          \codefile{evaluation/synthetic\_benchmark.py} so the benchmark
          can generate test variants.
\end{itemize}


% ─────────────────────────────────────────────────────────────────────
\section{Add an Entity Type}
\label{sec:add-entity-type}

\noindent\textbf{What you are adding.}
A new category of node in the knowledge graph (e.g.\ \ic{METRIC},
\ic{STANDARD}).

\medskip\noindent\textbf{Files to touch.}

\begin{enumerate}
    \item \codefile{apps/api/models/ontology.py} --- add to \ic{EntityType}
          enum.
    \item \codefile{apps/api/models/ontology.py} --- update
          \ic{VALID\_ENTITY\_RELATIONSHIPS} to declare which relationship
          types are legal for the new entity type.
    \item \codefile{apps/web/} --- frontend color map (entity badge
          colors).
    \item \codefile{apps/web/} --- graph node style (shape, size, color in
          the React Flow renderer).
\end{enumerate}

\medskip\noindent\textbf{Steps.}

\begin{enumerate}
    \item Add the new member to \ic{EntityType}:
\begin{lstlisting}[style=python]
class EntityType(Enum):
    TECHNOLOGY = "technology"
    CONCEPT = "concept"
    PATTERN = "pattern"
    SYSTEM = "system"
    PERSON = "person"
    ORGANIZATION = "organization"
    METRIC = "metric"  # <-- NEW
\end{lstlisting}

    \item For every relationship type in
          \ic{VALID\_ENTITY\_RELATIONSHIPS}, decide whether your new type
          participates.  Add the relevant \ic{(source, target)} tuples:

\begin{lstlisting}[style=python]
"RELATED_TO": {
    # ... existing pairs ...
    ("metric", "technology"),   # e.g., "latency RELATED_TO Redis"
    ("metric", "concept"),      # e.g., "throughput RELATED_TO caching"
},
\end{lstlisting}

    \item In the frontend, add a color entry for the new type so badges and
          graph nodes render correctly.  Search for the existing
          \ic{ENTITY\_TYPE\_COLORS} or equivalent mapping and add your type.
    \item Restart both backend and frontend, then ingest a conversation
          that mentions your new entity type.
\end{enumerate}

\medskip\noindent\textbf{Verification.}

\begin{lstlisting}[style=bash, numbers=none]
# Backend: run the test suite
cd apps/api && .venv/bin/pytest tests/ -v -k entity

# Frontend: check TypeScript compiles
cd apps/web && pnpm typecheck
\end{lstlisting}

\medskip\noindent\textbf{Gotchas.}

\begin{itemize}
    \item If you forget to update \ic{VALID\_ENTITY\_RELATIONSHIPS}, the
          validator (\codefile{services/validator.py}) will reject
          relationships involving the new type at runtime.
    \item The extraction prompt in \codefile{services/extractor.py}
          enumerates valid entity types.  Update the prompt so the LLM
          knows to produce the new type.
    \item Remember to bump \confkey{LLM\_EXTRACTION\_PROMPT\_VERSION} if
          you change the prompt (see Recipe~17.7).
\end{itemize}


% ─────────────────────────────────────────────────────────────────────
\section{Add a Relationship Type}
\label{sec:add-relationship-type}

\noindent\textbf{What you are adding.}
A new edge kind in the knowledge graph (e.g.\ \ic{REPLACES},
\ic{COMPLEMENTS}).

\medskip\noindent\textbf{Files to touch.}

\begin{enumerate}
    \item \codefile{apps/api/models/ontology.py} --- \ic{RelationType}
          enum, \ic{VALID\_ENTITY\_RELATIONSHIPS}, and the appropriate
          frozen set (\ic{ENTITY\_ONLY\_RELATIONSHIPS},
          \ic{DECISION\_ONLY\_RELATIONSHIPS}, or
          \ic{DECISION\_ENTITY\_RELATIONSHIPS}).
    \item \codefile{apps/api/services/extractor.py} --- detection logic in
          the extraction prompt.
    \item \codefile{apps/web/} --- edge style (color, dash pattern, label)
          in the graph renderer.
\end{enumerate}

\medskip\noindent\textbf{Steps.}

\begin{enumerate}
    \item Add the new member to \ic{RelationType}:

\begin{lstlisting}[style=python]
class RelationType(Enum):
    # ... existing members ...
    COMPLEMENTS = "COMPLEMENTS"  # X works well alongside Y
\end{lstlisting}

    \item Declare valid source/target type pairs:

\begin{lstlisting}[style=python]
VALID_ENTITY_RELATIONSHIPS["COMPLEMENTS"] = {
    ("technology", "technology"),  # Redis COMPLEMENTS PostgreSQL
    ("pattern", "pattern"),        # CQRS COMPLEMENTS Event Sourcing
}
\end{lstlisting}

    \item Add the new type to the correct frozen set:

\begin{lstlisting}[style=python]
ENTITY_ONLY_RELATIONSHIPS: frozenset[str] = frozenset([
    # ... existing ...
    "COMPLEMENTS",
])
\end{lstlisting}

    \item Update the extraction prompt in \codefile{services/extractor.py}
          to include the new relationship type and describe when the LLM
          should use it.
    \item Add an edge style in the frontend graph renderer (color, line
          dash, arrowhead).
    \item Bump \confkey{LLM\_EXTRACTION\_PROMPT\_VERSION}.
\end{enumerate}

\medskip\noindent\textbf{Verification.}

\begin{lstlisting}[style=bash, numbers=none]
cd apps/api
.venv/bin/pytest tests/ -v -k relationship
python -c "from models.ontology import validate_entity_relationship; \
  print(validate_entity_relationship('COMPLEMENTS','technology','technology'))"
\end{lstlisting}

The second command should print \ic{(True, None)}.

\medskip\noindent\textbf{Gotchas.}

\begin{itemize}
    \item \ic{ALL\_RELATIONSHIP\_TYPES} is a union of the three frozen
          sets.  It recomputes automatically, but only if you added your
          type to one of them.
    \item The \ic{validate\_entity\_relationship()} function normalizes
          inputs to uppercase/lowercase, so the enum value casing must be
          all-uppercase.
\end{itemize}


% ─────────────────────────────────────────────────────────────────────
\section{Add an API Endpoint}
\label{sec:add-api-endpoint}

\noindent\textbf{What you are adding.}
A new REST endpoint on the FastAPI backend.

\medskip\noindent\textbf{Files to touch.}

\begin{enumerate}
    \item \codefile{apps/api/routers/\textit{your\_module}.py} --- new
          router file.
    \item \codefile{apps/api/main.py} --- register the router with
          \ic{app.include\_router()}.
    \item Pydantic request/response models (inline or in
          \codefile{models/}).
\end{enumerate}

\medskip\noindent\textbf{Steps.}

\begin{enumerate}
    \item Create the router file:

\begin{lstlisting}[style=python, title=routers/health\_check.py]
"""Health-check router with database connectivity test."""

from fastapi import APIRouter, Depends
from pydantic import BaseModel

router = APIRouter()

class HealthResponse(BaseModel):
    status: str
    neo4j: bool
    postgres: bool

@router.get("/", response_model=HealthResponse)
async def health_check():
    """Return service health with dependency status."""
    # TODO: actual connectivity checks
    return HealthResponse(
        status="ok",
        neo4j=True,
        postgres=True,
    )
\end{lstlisting}

    \item Register the router in \codefile{main.py}:

\begin{lstlisting}[style=python]
from routers import health_check

app.include_router(
    health_check.router,
    prefix="/api/health",
    tags=["Health"],
)
\end{lstlisting}

    \item If the endpoint needs user context, add the \ic{user\_id}
          dependency that existing routers use for per-user data isolation.

    \item Write tests in \codefile{tests/routers/test\_health\_check.py}.
\end{enumerate}

\medskip\noindent\textbf{Verification.}

\begin{lstlisting}[style=bash, numbers=none]
# Start backend
cd apps/api && .venv/bin/uvicorn main:app --reload --port 8000

# In another terminal
curl -s http://localhost:8000/api/health/ | python -m json.tool
\end{lstlisting}

\medskip\noindent\textbf{Gotchas.}

\begin{itemize}
    \item Every data-returning endpoint should filter by \ic{user\_id} to
          maintain per-user isolation.  Review existing routers like
          \codefile{routers/decisions.py} for the pattern.
    \item Pydantic strict mode is used project-wide.  Use explicit type
          annotations; implicit coercion will raise validation errors.
    \item The trailing slash matters.  FastAPI redirects \ic{/api/health}
          to \ic{/api/health/} by default.
\end{itemize}


% ─────────────────────────────────────────────────────────────────────
\section{Add an LLM Provider}
\label{sec:add-llm-provider}

\noindent\textbf{What you are adding.}
A new LLM backend (e.g.\ Google Vertex AI, Ollama) alongside the existing
NVIDIA NIM and Amazon Bedrock providers.

\medskip\noindent\textbf{Files to touch.}

\begin{enumerate}
    \item \codefile{apps/api/services/llm\_providers/\textit{your\_provider}.py}
          --- implement \ic{BaseLLMProvider} (and optionally
          \ic{BaseEmbeddingProvider}).
    \item \codefile{apps/api/services/llm\_providers/\_\_init\_\_.py} ---
          register in the factory functions.
    \item \codefile{apps/api/config.py} --- add any new environment
          variables.
\end{enumerate}

\medskip\noindent\textbf{Steps.}

\begin{enumerate}
    \item Create your provider class by implementing the three abstract
          members of \ic{BaseLLMProvider}:

\begin{lstlisting}[style=python, title=llm\_providers/ollama.py]
"""Ollama LLM provider for local model inference."""

from typing import AsyncIterator
from services.llm_providers.base import BaseLLMProvider

class OllamaLLMProvider(BaseLLMProvider):
    async def generate(
        self,
        messages: list[dict],
        temperature: float = 0.6,
        max_tokens: int = 4096,
    ) -> tuple[str, dict]:
        # Call Ollama REST API
        ...

    async def generate_stream(
        self,
        messages: list[dict],
        temperature: float = 0.6,
        max_tokens: int = 4096,
    ) -> AsyncIterator[str]:
        # Stream from Ollama REST API
        ...

    @property
    def model_name(self) -> str:
        return "ollama/llama3"
\end{lstlisting}

    \item Register in the factory
          (\codefile{llm\_providers/\_\_init\_\_.py}):

\begin{lstlisting}[style=python]
def get_llm_provider() -> BaseLLMProvider:
    settings = get_settings()
    provider_name = getattr(settings, "llm_provider", "nvidia")

    if provider_name == "bedrock":
        from services.llm_providers.bedrock import BedrockLLMProvider
        return BedrockLLMProvider()
    elif provider_name == "ollama":           # <-- NEW
        from services.llm_providers.ollama import OllamaLLMProvider
        return OllamaLLMProvider()
    else:
        from services.llm_providers.nvidia import NvidiaLLMProvider
        return NvidiaLLMProvider()
\end{lstlisting}

    \item Add config variables to \codefile{config.py}:

\begin{lstlisting}[style=python]
# In Settings class
ollama_base_url: str = "http://localhost:11434"
ollama_model: str = "llama3"
\end{lstlisting}

    \item Set \confkey{LLM\_PROVIDER=ollama} in your \codefile{.env} file.
\end{enumerate}

\medskip\noindent\textbf{Verification.}

\begin{lstlisting}[style=bash, numbers=none]
LLM_PROVIDER=ollama python -c "
from services.llm_providers import get_llm_provider
p = get_llm_provider()
print(type(p).__name__, p.model_name)
"
\end{lstlisting}

\medskip\noindent\textbf{Gotchas.}

\begin{itemize}
    \item The \ic{generate()} method must return a \ic{(text, usage\_dict)}
          tuple.  The \ic{usage\_dict} should contain
          \ic{prompt\_tokens}, \ic{completion\_tokens}, and
          \ic{total\_tokens} for token logging.
    \item If your provider does not support streaming, raise
          \ic{NotImplementedError} from \ic{generate\_stream()} rather
          than silently failing.
    \item The \ic{LLMClient} in \codefile{services/llm.py} wraps
          providers with rate limiting and retry logic.  You do not need to
          implement retries inside your provider.
\end{itemize}


% ─────────────────────────────────────────────────────────────────────
\section{Add an Evaluation Script}
\label{sec:add-evaluation-script}

\noindent\textbf{What you are adding.}
A new evaluation or analysis script that follows the project's conventions.

\medskip\noindent\textbf{Files to touch.}

\begin{enumerate}
    \item \codefile{apps/api/evaluation/\textit{your\_script}.py} --- new
          script.
\end{enumerate}

\medskip\noindent\textbf{Steps.}

\begin{enumerate}
    \item Follow the established pattern: argparse for CLI arguments, JSON
          output for machine-readable results, and use synthetic
          conversations as input data.

\begin{lstlisting}[style=python, title=evaluation/my\_baseline.py]
"""TF-IDF cosine similarity baseline for entity resolution.

Usage:
    python -m evaluation.my_baseline [--output PATH]
"""

import argparse
import json
from pathlib import Path

def main():
    parser = argparse.ArgumentParser(
        description="TF-IDF baseline evaluation"
    )
    parser.add_argument(
        "--output",
        type=Path,
        default=Path("evaluation/data/v5/tfidf_results.json"),
    )
    args = parser.parse_args()

    # Load synthetic conversations
    conv_dir = Path("evaluation/data/synthetic_conversations")
    conversations = []
    for f in sorted(conv_dir.glob("*.json")):
        conversations.append(json.loads(f.read_text()))

    # Run evaluation ...
    results = {"precision": 0.0, "recall": 0.0, "f1": 0.0}

    # Write results
    args.output.write_text(
        json.dumps(results, indent=2)
    )
    print(f"Results written to {args.output}")

if __name__ == "__main__":
    main()
\end{lstlisting}

    \item Use the same train/test split as the benchmark (140/60) to ensure
          comparable results.
    \item Output JSON with at least \ic{precision}, \ic{recall}, and
          \ic{f1} keys.
\end{enumerate}

\medskip\noindent\textbf{Verification.}

\begin{lstlisting}[style=bash, numbers=none]
cd apps/api
python -m evaluation.my_baseline --output /tmp/test_results.json
cat /tmp/test_results.json
\end{lstlisting}

\medskip\noindent\textbf{Gotchas.}

\begin{itemize}
    \item Always run from \codefile{apps/api/} so that relative imports
          resolve correctly (\ic{sys.path} manipulation at the top of each
          evaluation script handles this).
    \item Write results to \codefile{evaluation/data/v5/} for consistency
          with existing baselines.
    \item If your script calls the LLM, use
          \ic{sanitize\_input=False} when passing synthetic conversations
          to avoid the prompt sanitizer blocking ``USER:''/``ASSISTANT:''
          labels.
\end{itemize}


% ─────────────────────────────────────────────────────────────────────
\section{Modify the Extraction Prompt}
\label{sec:modify-extraction-prompt}

\noindent\textbf{What you are adding.}
Changes to the LLM prompt template that governs decision extraction from
conversations.

\medskip\noindent\textbf{Files to touch.}

\begin{enumerate}
    \item \codefile{apps/api/services/extractor.py} --- the prompt
          template string.
    \item \codefile{apps/api/config.py} --- the
          \confkey{llm\_extraction\_prompt\_version} setting.
\end{enumerate}

\medskip\noindent\textbf{Steps.}

\begin{enumerate}
    \item Edit the prompt in \codefile{services/extractor.py}.
    \item Open \codefile{config.py} and bump
          \confkey{llm\_extraction\_prompt\_version} (e.g.\ from
          \ic{"v1"} to \ic{"v2"}).
    \item Re-run extraction on a small sample to verify the new prompt
          produces valid JSON output.
\end{enumerate}

\begin{warning}
Always bump the prompt version after changing the prompt template.
\sysname{} uses a Redis-backed LLM response cache keyed partly by prompt
version.  If you change the prompt without bumping the version, the system
will continue serving \emph{old cached responses} generated by the
previous prompt.  This is a silent correctness bug that is difficult to
diagnose.
\end{warning}

\medskip\noindent\textbf{Verification.}

\begin{lstlisting}[style=bash, numbers=none]
cd apps/api

# Test on a single conversation
python -c "
import asyncio, json
from services.extractor import DecisionExtractor

async def test():
    ext = DecisionExtractor()
    result = await ext.extract_from_conversation(
        json.load(open(
            'evaluation/data/synthetic_conversations/conv-001.json'
        )),
        sanitize_input=False,
    )
    print(json.dumps(result, indent=2, default=str))

asyncio.run(test())
"
\end{lstlisting}

\medskip\noindent\textbf{Gotchas.}

\begin{itemize}
    \item NVIDIA Nemotron generates \ic{<think>...</think>} blocks that
          consume tokens.  Set \ic{max\_tokens=4096} or higher to avoid
          truncated JSON output.
    \item The \ic{strip\_thinking\_tags()} function in
          \codefile{services/llm.py} removes thinking blocks
          automatically, but only \emph{after} the full response is
          received.
    \item After a prompt change, run the full evaluation suite
          (\codefile{evaluation/run\_all.py}) to measure impact on
          B-cubed metrics before merging.
\end{itemize}
